\vspace{-50pt}


\newpage
\section{Introduction}
Proteins perform a vast array of functions within organisms. As the evolution process spans millions to billions of years, the field of protein design aims to automate and accelerate it through computational methods. Since the function of proteins is confered via the structure, which is determined by the sequence, a line of research focuses on generating sequences \citep{dauparasRobustDeepLearning2022, zhengStructureinformedLanguageModels2023, alamdariProteinGenerationEvolutionary2024, wang2024dplm} and structures \cite{yimSE3DiffusionModel2023b, watsonNovoDesignProtein2023, ingrahamIlluminatingProteinSpace2023, zhang2025odesign}. Recently, a new paradigm has emerged emphasizing the co-generation of sequence and structure \cite{wang2025dplm2, hsieh2025dplm2_1, chen2025apm, multiflow}. However, improving sample diversity while preserving the coherence between modalities, specifically sequence-structure consistency, remains a formidable challenge. To address this, recent works have pivoted toward all-atom modeling \cite{quPallatomUnlockingNew2025, geffnerLaProteinaAtomisticProtein2025}. By directly designing proteins at the all-atom level, these methods naturally derive amino acid identities from atomic coordinates, achieving superior consistency. These methodological advancements have significantly propelled the field of \textit{de novo} protein design, particularly in binder design, the creation of mini-binders for specific therapeutic targets \cite{zambaldi2024novo}. This domain has witnessed a rapid evolution, progressing from classical hallucination-based methods \cite{pacesaOneshotDesignFunctional2025} or backbone-based design model \cite{zhang2025odesign, ren2025pxdesign} to the current wave of all-atom diffusion models \cite{butcher2025novo, chai2025zero, team2025latent}, all striving for enhanced designability and binding affinity.


In this study, we introduce \texttt{SeedProteo}, a diffusion-based model that learns the manifold of proteins at the all-atom level. SeedProteo mimics the architecture of AlphaFold3 \cite{abramsonAccurateStructurePrediction2024, wohlwendBoltz1DemocratizingBiomolecular2024, passaroAccurateEfficientBinding}, which has proven its effectiveness in protein folding. Through unifying the representation of different amino acids, we significantly enhance its generative power when the protein sequence and homology are not provided. A major issue in all-atom modeling is the backbone-sidechain inconsistency (i.e., the sequence-structure inconsistency), which means that the sidechain atoms are locally plausible but the derived sequence fails to fold into the global backbone structure. SeedProteo iteratively evolves the sequence and structure and integrates two components to enhance the modeling ability of two modalities:
(i) SeedProteo enhances the structural feasibility through the integration of geometric reasoning in the denoising process. To be specific, we reuse the denoised structure from the previous timestep \cite{chen2022analog} to stabilize the sampling process and conduct geometric reasoning in the pair representation space.
(ii) SeedProteo conducts energy minimization within the sequence space modeled by a Markov random field (MRF) \cite{metzler2005markov}. By integrating high-order couplings among distinct positions, we are able to sample globally feasible sequences and prevent being trapped in the local minima of simple patterns constituted by high-frequency amino acids.
We conduct extensive experiments on unconditional sampling, binder design, and in vitro validation. The results demonstrate that:

1) In unconditional generation, SeedProteo exhibits superior length generalization and structural diversity compared to state-of-the-art baselines. It maintains stable performance on long sequences (up to 1000 residues) and complex topologies, effectively addressing the length-dependent degradation issues observed in baseline diffusion models.


2) In binder design, to address the trade-off between design success rate and structural diversity, we introduce two inference modes: \texttt{SeedProteo-R} (Robust), which prioritizes longer secondary structure segments for higher design success, and \texttt{SeedProteo-D} (Diverse), which fosters topological richness. Experiments show robust performance across diverse targets, particularly excelling in challenging cases.


3) Furthermore, our analysis offers specific insights into how to achieve superior sequence-structure consistency within the \texttt{atom14} representation schema \cite{quPallatomUnlockingNew2025, butcher2025novo, starkUniversalBinderDesign}. We demonstrate that introducing a global MRF-based decoding module leads to more self-consistent outputs.


4) Beyond in-silico evaluation, we further validate SeedProteo with wet-lab experiments, demonstrating that designed binders exhibit strong binding in SPR assays with high hit rates across multiple therapeutic targets.

